\chapter{Pembagian dan Kelipatan}\label{ch:g5:division}

Pembagian bersusun yang dipelajari pada \cref{ch:g4:division}
mendapat dua peningkatan: \omterm{def:g5:division:multiple}{pembagi} dua angka, dan \omterm{def:g3:division:remainder}{hasil bagi} yang
berlanjut melewati titik --- $3 \div 4$ akhirnya punya jawaban,
$0.75$. Bab ini ditutup dengan \omterm{def:g5:division:multiple}{kelipatan}, \omterm{def:g5:division:multiple}{pembagi} dan jalan pintas
keterbagian yang pertama.

\section{Membagi dengan bilangan dua angka}

\begin{method}[Pembagi dua angka]\label{met:g5:division:twodigits}
Caranya sama seperti sebelumnya, dengan satu keterampilan tambahan:
menerka berapa kali \omterm{def:g5:division:multiple}{pembaginya} muat. Tulislah dulu tabel kecil
\omterm{def:g5:division:multiple}{kelipatan} \omterm{def:g5:division:multiple}{pembaginya}. Untuk $986 \div 23$ (\omterm{def:g5:division:multiple}{kelipatan} 23: 23, 46, 69,
92, 115, 138, 161, 184, 207):
\begin{enumerate}
  \item $98$ puluhan: $23 \times 4 = 92$ muat, $23 \times 5 = 115$
        tidak: angkanya $4$, \omterm{def:g3:division:remainder}{sisanya} $6$;
  \item turunkan angka $6$: menjadi $66$: $23 \times 2 = 46$ muat:
        angkanya $2$, \omterm{def:g3:division:remainder}{sisanya} $20$;
  \item \omterm{def:g3:division:remainder}{hasil baginya} $42$, \omterm{def:g3:division:remainder}{sisanya} $20$. Periksa:
        $23 \times 42 + 20 = 966 + 20 = 986$.
\end{enumerate}
\end{method}

\section{Hasil bagi desimal}

\begin{example}[Pembagian yang menolak berhenti pada sisa]\label{ex:g5:division:decimal}
Bagilah $3$ pizza kepada $4$ anak: $3 \div 4$ punya \omterm{def:g3:division:remainder}{hasil bagi} $0$
dan \omterm{def:g3:division:remainder}{sisa} $3$ --- tidak berguna! Lanjutkan pembagiannya melewati titik:
\begin{enumerate}
  \item $3$ satuan $= 30$ persepuluhan; $30 \div 4$: $7$
        persepuluhan, \omterm{def:g3:division:remainder}{sisanya} $2$ persepuluhan;
  \item $2$ persepuluhan $= 20$ perseratusan; $20 \div 4 = 5$
        perseratusan, \omterm{def:g3:division:remainder}{sisanya} $0$.
\end{enumerate}
Jadi $3 \div 4 = 0.75$: setiap anak mendapat $0.75$ pizza --- tiga
perempat, seperti yang sudah diketahui oleh bahasa \omterm{def:g4:fractions:def}{pecahan}
(\cref{prop:g5:fractions:decimal}).
\end{example}

\begin{method}[Melanjutkan pembagian melewati titik]\label{met:g5:division:continue}
Ketika satuannya habis dan masih ada \omterm{def:g3:division:remainder}{sisa}:
\begin{enumerate}
  \item tulislah \omterm{def:g4:decimals:point}{titik desimal} pada \omterm{def:g3:division:remainder}{hasil baginya};
  \item tambahkan sebuah \omterm{def:g1:counting:zero}{nol} pada \omterm{def:g3:division:remainder}{sisanya} (satuan menjadi
        persepuluhan, persepuluhan menjadi perseratusan\,\dots) lalu
        teruskan membagi;
  \item berhentilah ketika \omterm{def:g3:division:remainder}{sisanya} $0$ --- atau ketika angkanya
        mulai berulang selamanya, seperti $1 \div 3 = 0.333\dots$:
        lalu berikan nilai yang \omterm{def:g4:numbers:round}{dibulatkan}.
\end{enumerate}
\end{method}

\begin{example}\label{ex:g5:division:continue}
$27 \div 6$: \omterm{def:g3:division:remainder}{hasil bagi} $4$, \omterm{def:g3:division:remainder}{sisa} $3$; pasang titiknya; $30$
persepuluhan $\div\ 6 = 5$: jadi $27 \div 6 = 4.5$. \quad
$22 \div 8$: $2$, \omterm{def:g3:division:remainder}{sisa} $6$; lalu $60 \div 8 = 7$ \omterm{def:g3:division:remainder}{sisa} $4$;
lalu $40 \div 8 = 5$: jadi $22 \div 8 = 2.75$.
\end{example}

\section{Kelipatan dan pembagi}

\begin{definition}[Kelipatan, pembagi]\label{def:g5:division:multiple}
\emph{Kelipatan}\index{kelipatan} $6$ adalah tabelnya yang
dilanjutkan selamanya: $0, 6, 12, 18, 24, \dots$ Ketika $24$
merupakan kelipatan $6$, kita juga mengatakan bahwa $6$ adalah
\emph{pembagi}\index{pembagi} dari $24$, atau bahwa $24$
\emph{habis dibagi} $6$: pembagian $24 \div 6$ tidak menyisakan apa-apa.
\end{definition}

\begin{proposition}[Jalan pintas keterbagian]\label{prop:g5:division:criteria}
Tanpa membagi, sebuah bilangan habis dibagi:
\begin{itemize}
  \item oleh $2$ jika angka terakhirnya $0$, $2$, $4$, $6$ atau $8$
        (yaitu \omterm{def:g3:numbers:evenodd}{bilangan genap});
  \item oleh $5$ jika angka terakhirnya $0$ atau $5$;
  \item oleh $10$ jika angka terakhirnya $0$;
  \item oleh $3$ jika \emph{\omterm{def:g1:addition:def}{jumlah} angka-angkanya} habis dibagi
        $3$ (misalnya $741$: $7 + 4 + 1 = 12$, habis dibagi --- dan
        memang $741 = 3 \times 247$).
\end{itemize}
\end{proposition}
\admitted

\begin{omfigure}
\begin{tikzpicture}[scale=0.5]
  \foreach \n in {0,...,29} {
    \pgfmathtruncatemacro{\r}{int(\n/10)}
    \pgfmathtruncatemacro{\c}{mod(\n,10)}
    \pgfmathtruncatemacro{\lbl}{\n+1}
    \pgfmathtruncatemacro{\mtwo}{mod(\lbl,2)}
    \pgfmathtruncatemacro{\mthree}{mod(\lbl,3)}
    \ifnum\mtwo=0
      \fill[omDef!25] (\c,-\r) ++(-0.45,-0.45) rectangle ++(0.9,0.9);
    \fi
    \ifnum\mthree=0
      \draw[omThm, thick] (\c,-\r) circle (0.42);
    \fi
    \node[font=\footnotesize] at (\c,-\r) {\lbl};
  }
  \node[font=\small, below] at (4.5,-2.8)
    {persegi biru: kelipatan $2$;\ \ lingkaran merah: kelipatan
    $3$;\ \ keduanya: kelipatan $6$};
\end{tikzpicture}

{\small Bilangan $1$ sampai $30$: \omterm{def:g5:division:multiple}{kelipatan} $2$ dan \omterm{def:g5:division:multiple}{kelipatan} $3$
membentuk pola --- dan bilangan yang memakai kedua tanda ($6$, $12$,
$18$, $24$, $30$) tepatnya adalah \omterm{def:g5:division:multiple}{kelipatan} $6$.}
\end{omfigure}

\begin{example}\label{ex:g5:division:criteria}
Apakah $534$ habis dibagi $2$? Berakhir dengan $4$: ya. Oleh $5$?
Tidak. Oleh $3$? $5 + 3 + 4 = 12$: ya. Jadi $534$ juga habis dibagi
$6$ (oleh $2$ \emph{dan} oleh $3$) --- periksa: $534 = 6 \times 89$.
\end{example}

\section{Latihan}

\begin{exercise}[$\star$]\label{exo:g5:division:1}
Hitunglah pembagian bersusunnya (tulis tabel \omterm{def:g5:division:multiple}{kelipatannya} dulu):
$851 \div 23$; \quad $704 \div 32$.
\end{exercise}

\begin{exercise}[$\star$]\label{exo:g5:division:2}
Hitunglah \omterm{def:g3:division:remainder}{hasil bagi} desimal yang tepat: $9 \div 2$; \quad
$27 \div 4$; \quad $33 \div 5$; \quad $21 \div 8$.
\end{exercise}

\begin{exercise}[$\star$]\label{exo:g5:division:3}
Empat orang teman membagi rata tagihan di restoran sebesar $54$. Berapa
yang dibayar masing-masing? (Lanjutkan melewati titik.)
\end{exercise}

\begin{exercise}[$\star$]\label{exo:g5:division:4}
Hitunglah $10 \div 3$, dilanjutkan tiga angka melewati titik. Apa
yang kamu perhatikan? Berikan nilainya yang \omterm{def:g4:numbers:round}{dibulatkan} ke perseratusan.
\end{exercise}

\begin{exercise}[$\star$]\label{exo:g5:division:5}
Sebutkan: \omterm{def:g5:division:multiple}{kelipatan} $7$ sampai $70$; \omterm{def:g5:division:multiple}{pembagi} dari $18$; dan \omterm{def:g5:division:multiple}{pembagi}
dari $24$.
\end{exercise}

\begin{exercise}[$\star$]\label{exo:g5:division:6}
Habis dibagi $2$? oleh $5$? oleh $10$? oleh $3$? Ujilah setiap jalan
pintas pada: $470$; \quad $735$; \quad $8\,001$; \quad $1\,314$.
\end{exercise}

\begin{exercise}[$\star$]\label{exo:g5:division:7}
Carilah semua bilangan antara $60$ dan $80$ yang habis dibagi $3$
\emph{dan} oleh $5$. (Keterbagian tunggal apa yang digabungkan oleh
keduanya?)
\end{exercise}

\begin{exercise}[$\star$]\label{exo:g5:division:8}
Sehelai pita sepanjang $7.2$ m dipotong menjadi $6$ bagian yang sama
panjang. Berapa panjang tiap bagian? (Bagilah bilangan desimal
dengan bilangan bulat: bagilah meternya dulu, lalu persepuluhannya.)
\end{exercise}

\begin{exercise}[$\star$]\label{exo:g5:division:9}
$1\,000$ kelereng dikemas ke dalam kantong berisi $24$. Berapa
kantong yang penuh, dan berapa kelereng yang tersisa? (\omterm{def:g5:division:multiple}{Pembagi} dua angka.)
\end{exercise}

\begin{exercise}[$\star\star$]\label{exo:g5:division:10}
Dapatkah $5$ orang teman membagi rata $7$ batang cokelat itu tanpa
\omterm{def:g3:division:remainder}{sisa}?
Berikan bagian masing-masing sebagai bilangan desimal. Pertanyaan
yang sama untuk $7$ teman yang membagi $5$ batang --- apa yang berbeda?
\end{exercise}

\begin{exercise}[$\star\star$]\label{exo:g5:division:11}
Dengan jalan pintas \omterm{def:g1:addition:def}{jumlah} angka, carilah angka $d$ terkecil yang
membuat $52d$ (bilangan tiga angka yang berakhir dengan $d$) habis
dibagi $3$. Lalu carilah semua angka $d$ yang demikian.

\end{exercise}
